\section{Introduction and theorem}

Egyptian-fraction problems ask for representations of rational numbers as sums
of reciprocals subject to arithmetic restrictions on the denominators.  Erd\H{o}s
and Graham asked whether every positive rational with squarefree reduced
denominator admits such a representation using distinct products of two primes
\cite{ErdosGraham1980}; the current problem record is also maintained by Bloom
\cite{BloomErdos306}.  We prove that the evident squarefreeness obstruction is
the only one.

An integer is a \emph{squarefree semiprime} if it is a product $pq$ of two
distinct primes.

\begin{lemma}[Semiprime criterion]\label{lem:semiprime-equivalence}
For $n\ge1$, the following are equivalent:
\begin{enumerate}[label=(\roman*)]
\item $n=pq$ for two distinct primes $p,q$;
\item $\omega(n)=\Omega(n)=2$.
\end{enumerate}
\end{lemma}
\begin{proof}
If $n=pq$ with $p\ne q$, its prime factorization has two factors counted with
multiplicity and two distinct prime factors.  Conversely, $\Omega(n)=2$ gives
$n=pq$ for primes $p,q$, possibly equal.  Equality would imply $\omega(n)=1$,
so $p\ne q$.
\end{proof}

\begin{theorem}[Erd\H{o}s Problem 306]\label{thm:headline}
Let $q>0$ be rational and write $q=a/b$ in lowest terms.  Then $q$ is a finite
sum of reciprocals of distinct squarefree semiprimes if and only if $b$ is
squarefree.
\end{theorem}

The necessity is immediate.  If $S$ is a finite set of squarefree integers and
$M=\lcmop(S)$, then
\[
  \sum_{n\in S}\frac1n
  =\frac{\sum_{n\in S}M/n}{M}.
\]
The reduced denominator divides $M$, and an lcm of squarefree integers is
squarefree.  Every divisor of a squarefree integer is squarefree, so cancellation
causes no difficulty.

For the converse we retain an avoidance parameter.  For a finite set
$T\subset\N$, write $\operatorname{Rep}_T(r)$ if there is a finite set $S$ of
squarefree semiprimes such that
\[
 S\cap T=\varnothing,
 \qquad
 \sum_{n\in S}\frac1n=r.
\]
The obstruction set is part of the mechanism: it prevents denominators from
being repeated when representations are added.

\begin{lemma}[Numerator reduction]\label{lem:numerator-reduction}
If $\operatorname{Rep}_T(1/b)$ holds for every finite $T$, then
$\operatorname{Rep}_T(a/b)$ holds for every integer $a\ge1$ and every finite
$T$.
\end{lemma}
\begin{proof}
Induct on $a$.  Given a set $S_a$ representing $a/b$, apply the unit-numerator
assertion with obstruction set $T\cup S_a$.  The new set is disjoint from
$S_a$, and the reciprocal sums add.
\end{proof}

\begin{lemma}[Small denominators]\label{lem:small-b}
If the avoiding unit-numerator assertion holds for every squarefree $b\ge3$,
then it holds for all positive squarefree $b$.
\end{lemma}
\begin{proof}
For $b=2$, take disjoint avoiding representations of $1/3$ and $1/6$ and use
$1/2=1/3+1/6$.  For $b=1$, take pairwise disjoint avoiding representations of
$1/2$, $1/3$, and $1/6$, and use $1=1/2+1/3+1/6$.
\end{proof}

It remains to prove the following avoiding statement.

\begin{theorem}[Avoiding unit-fraction construction]\label{thm:structural-construction}
Let $b\ge3$ be squarefree and let $T\subset\N$ be finite.  Then
$\operatorname{Rep}_T(1/b)$ holds.
\end{theorem}

The proof is organized by mathematical handoffs rather than by a numerical
certificate.  First, one common analytic threshold supplies dense dyadic prime
blocks and reciprocal-prime mass greater than one.  Finite Fourier inversion
then reduces the construction to a strict major/minor comparison.  A centred
CRT control graph yields local rigidity, a corrected adjacent-label theorem,
and a global localization estimate.  A pair family completes the reciprocal
load and induces a compact interval of Bernoulli weights.  The major frequencies
have a Gaussian lower bound, while two exact minor classes are controlled by
localization and by a common reservoir of auxiliary primes.  Three symbolic
budget shares close the comparison; every parameter is fixed before the bottom
dyadic scale.
